\section{Finite-count protection and higher integrability of the raw density}
\label{sec:v54-density-integrability}
We return to the scalar physical density. A fixed-band positive upper
bound is less precise than the local limit theorem, but it has a useful
advantage here: its error can be retained uniformly as the protection
scale decreases. We use this bound to obtain a polynomial height budget
for the protected source. The thin-layer estimates then control the
mass above a large density level. The complete exponential finite-count
height bound closes the remaining tail integration. None of these
steps continues a trajectory across a physical singularity.

Use the original exact labels and normalization
\[
 P(u)=m^2p_{n,R}(k,m,u),\qquad
 G(u)=\mathcal L_{m,R}(k_1,k_2,u,n),\qquad 1\le n\le m.
\]
The section source is $\nu_R^*=\nu|_{Y_R^*}/c$. In particular $P$
is nonnegative, whereas $G$ is the real arithmetic transition kernel
and need not be positive at every finite count. All density statements
are almost everywhere. The constants below are uniform in $R,n,k$
and in translated roof intervals.

\subsection{One auxiliary band for all protection scales}
Fix once an auxiliary roof band $B_\circ\ge B_0$. It is independent
of the collision count, the protection scale, the roof interval and
the density level to be introduced below.

\begin{lemma}[Finite-count positive window bounds at one band]
\label{lem:v54-fixed-band-windows}
There are $C<\infty$ and $0<\rho<1$ such that, for every interval
$I$ of length $h>0$ and $m\ge2$,
\begin{equation}\label{eq:v54-unmarked-window}
 m^2\nu_R^*\{K_{n,R}=k,N_{n,R}=m,T_{n,R}\in I\}
       \le C(1+h)(1+m^2\rho^{m/2}).
\end{equation}
For the full guard $\Xi_{m,R}^{\varepsilon}$ in
\eqref{eq:v46-source-guard}, define the positive densities
\begin{equation}\label{eq:v54-positive-guard-split}
 F_\varepsilon=m^2f^{\varepsilon,1}_{n,k,m,R},\qquad
 R_\varepsilon=P-F_\varepsilon\ge0.
\end{equation}
Then, simultaneously for all $0<\varepsilon<\varepsilon_0$,
\begin{equation}\label{eq:v54-all-margin-window}
 \int_I R_\varepsilon(u)\,du
 \le C(1+h)\bigl(\varepsilon^{1/16}+m^3\rho^{m/2}\bigr).
\end{equation}
The scale $\varepsilon$ may depend on $m$ in these finite inequalities.
\end{lemma}
\begin{proof}
On the full collision source, the unmarked characteristic pairing is
$\ell(\mathscr T_{R,\theta}^m\nu)$. The fixed-band moving-peak
expansion used in Theorem~\ref{thm:v49-physical-marked-local}, with
no intermediate multiplier, bounds its four-frequency integral by
$C m^{-2}+C\rho^{m/2}$. Indeed, on each peak the amplitude is bounded
and the real quadratic damping has a common positive lower bound;
the Gaussian integral is $O(m^{-2})$. The fixed-band complementary
powers are exponential. The finite number of moving peak charts is
uniform, and no occupation resonance is removed.

Use a positive interval majorant of mass $h+2\pi/B_\circ$ and
Fourier support in this one fixed band. Its Fourier supremum is at
most its mass. Torus orthogonality and roof inversion therefore give
$C(h+B_\circ^{-1})+C(1+B_\circ h)m^2\rho^{m/2}$ for the
normalized positive probability. Dropping initial and terminal section
membership on the upper side and using $\nu_R^*\le\nu/c$ proves
\eqref{eq:v54-unmarked-window}. Constants depending on the fixed
$B_\circ$ have now been absorbed into $C$. The argument is finite
in $m$ and uniform in the interval length; it is not an application
of a fixed-window limit to a shrinking interval.

Write $q_0=47/53$, $d_j=\min(j,m-j)$ and
$d'_j=\min(j,m-1-j)$. Since the factors of $\Xi$ lie in $[0,1]$,
$1-\Xi$ is bounded by the union sum of
\[
 c_j<2\varepsilon q_0^{d_j/4},\qquad
 s_j<2\varepsilon q_0^{d_j/4},\qquad
 \Delta_j<2\varepsilon q_0^{d'_j/4}.
\]
There are $3m+2$ terms before overlaps are removed. Section strips
are bounded by Theorem~\ref{thm:v48-marked-strip-local}; incidence
and clearance strips are bounded by
Theorem~\ref{thm:v49-physical-marked-local}. All use $B_\circ$.
The clearance of flight $j$ is read at collision $j+1$. The mark
indices range from zero to $m$, while occupation still sums exactly
the times $0,\ldots,m-1$.

At most two contacts or flights have any prescribed depth, and
$\sum_{a\ge0}q_0^{a/64}<\infty$. Thus the leading thin-layer terms
sum to $C(1+h)\varepsilon^{1/16}$ and the width-independent errors
to $C(1+h)m^3\rho^{m/2}$. The same enlargement by $1/c$ proves
\eqref{eq:v54-all-margin-window} on the original exact-label source.
Take the largest of the finitely many spectral radii used above.
The estimates being summed are uniform in the widths and marks,
which proves the simultaneous assertion.
\end{proof}

\begin{proposition}[A polynomial protection loss in height]
\label{prop:v54-protected-height}
There is a constant $C_A$ such that for every $m\ge2$ and
$0<\varepsilon<\varepsilon_0$,
\begin{equation}\label{eq:v54-protected-height}
 \sup_{R,n,k}\|F_\varepsilon\|_\infty
 \le C_A\varepsilon^{-9}(1+m^2\rho^{m/2})
 \le C_A'\varepsilon^{-9}.
\end{equation}
The constants do not depend on $m$ or $\varepsilon$. In particular
this is a finite-count assertion, rather than a separate limsup
for each fixed protection scale.
\end{proposition}
\begin{proof}
Choose an absolute $d>0$, sufficiently small for
Lemma~\ref{lem:v46-near-critical-completion}, and put
\[
 \delta=d\varepsilon^3,\qquad
 h_\varepsilon'=c_1\varepsilon^3\delta^2=c_1d^2\varepsilon^9,
 \qquad r_\varepsilon=C_*\delta^2=C_*d^2\varepsilon^6.
\]
Reduce $d$ so that $r_\varepsilon<h_{\varepsilon/2}$ and both roof
widths are at most one. The guard is supported in
$\mathsf P_{m,R}(\varepsilon)$.

On $|\nabla F|\ge\delta$, the positive flow-box comparison
\eqref{eq:v46-slice-window} bounds the density by the original
exact-label probability in a window of length $2h_\varepsilon'$,
divided by $2h_\varepsilon'/e$. Apply the finite inequality
\eqref{eq:v54-unmarked-window}, not
\eqref{eq:v46-window-budget-limit}. The normalized height is at
most $C\varepsilon^{-9}(1+m^2\rho^{m/2})$. This use of a short
window pays its inverse length explicitly; it does not enlarge the
auxiliary Fourier band.

For the part $|\nabla F|<\delta$, complete each point to its
$\varepsilon/2$-protected normal critical center using
Lemma~\ref{lem:v46-near-critical-completion}. If the point has
roof $u$, its critical roof lies in $[u-r_\varepsilon,u]$.
The mass collars of width $r_\varepsilon$ have density at most
$A_*J_z$ and mass at least $a_*r_\varepsilon J_z$, by
Proposition~\ref{prop:v45-mass-collar}. They have the same original
$(n,k,m)$ labels. Different physical collars are disjoint and a word
has at most one center. Consequently
\[
 m^2\sum_{t_z\in[u-r_\varepsilon,u]}J_z
 \le \frac{C m^2}{r_\varepsilon}
 \nu_R^*\{K_{n,R}=k,N_{n,R}=m,
                         |T_{n,R}-u|<r_\varepsilon\}.
\]
No normal-strip smallness is needed for this upper bound; removing
that condition only enlarges the positive comparison. The collar
density upper bound and \eqref{eq:v54-unmarked-window} now give
$C\varepsilon^{-6}(1+m^2\rho^{m/2})$ for the normalized low-gradient
height. It is smaller than the required budget. Every collar is
charged once and the source is still $\nu_R^*$, not a renormalized
collar law. Add the two contributions. Finally
$\sup_{m\ge2}m^2\rho^{m/2}<\infty$ proves the second inequality.
\end{proof}

\subsection{A finite-error high-density argument}
The following measure-theoretic lemma states exactly how the exponentially
small finite-count error is used. It avoids interchanging a limsup with
an integral over unbounded density levels.

\begin{lemma}[Protection, thin mass and an exponential cap]
\label{lem:v54-tail-principle}
Let $P_m\ge0$ on a finite measure space. Suppose that, for constants
$K,\alpha,\kappa,\Gamma>0$ and every $0<\varepsilon<\varepsilon_0$,
\[
 P_m=A_{m,\varepsilon}+B_{m,\varepsilon},\quad
 A_{m,\varepsilon},B_{m,\varepsilon}\ge0,\quad
 \|A_{m,\varepsilon}\|_\infty\le C_A\varepsilon^{-K},
 \quad \int B_{m,\varepsilon}\le C_I(\varepsilon^\alpha+e^{-\kappa m}),
\]
with the same constants for all $m,\varepsilon$. Assume also
$\|P_m\|_\infty\le C_H e^{\Gamma m}$ and a uniform bound on
$\int P_m$. There are constants $L_0,C_I'$ such that for $L\ge L_0$,
\begin{equation}\label{eq:v54-tail-principle}
 \int_{\{P_m>L\}}P_m
      \le C_I'(L^{-\alpha/K}+e^{-\kappa m}).
\end{equation}
For every $0<s<\min\{\alpha/K,\kappa/\Gamma\}$ one has
$\sup_m\int P_m^{1+s}<\infty$.
\end{lemma}
\begin{proof}
Take $\varepsilon=(2C_A/L)^{1/K}$ and increase $L_0$ so that this
scale is admissible. Then $A_{m,\varepsilon}\le L/2$. On
$\{P_m>L\}$ this implies $B_{m,\varepsilon}\ge P_m/2$, which
proves \eqref{eq:v54-tail-principle}. The decomposition is exact at
this chosen scale for each finite $m$.

Let $H_m=\max\{1,C_H e^{\Gamma m}\}$. The layer-cake identity,
with an upper bound on the part $P_m\le1$, gives
\[
 \int P_m^{1+s}
 \le \int P_m+s\int_1^{H_m}L^{s-1}
                    \int_{\{P_m>L\}}P_m\,dL.
\]
The bounded interval up to $L_0$ is harmless. The algebraic tail
integrates uniformly when $s<\alpha/K$. The finite-error term is
at most $C_I'e^{-\kappa m}H_m^s$, which is bounded, and tends to
zero, when $s\Gamma<\kappa$. This integrates the finite error
before taking a collision limit. A collection of separate asymptotic
tail bounds, without this error and the cap, would not justify the
same conclusion.
\end{proof}

Choose constants once as follows. With $\rho$ from
Lemma~\ref{lem:v54-fixed-band-windows}, put
\[
 \kappa=-\tfrac14\log\rho>0.
\]
The polynomial errors are bounded by $C e^{-\kappa m}$ after
changing $C$. If $A>1$ is the complete finite-count height constant
in \eqref{eq:v44-full-roof-height}, put $\Gamma=\log(2A)$.
Equation~\eqref{eq:v44-weighted-height} then gives
$\|P\|_\infty\le C_H e^{\Gamma m}$. Define
\begin{equation}\label{eq:v54-admissible-exponent}
 s_* =\frac12\min\left\{\frac1{144},\frac{\kappa}{\Gamma}\right\}>0,
 \qquad q_*=1+s_*.
\end{equation}
All constants arise from one fixed auxiliary band and the complete
finite-count geometry. No numerical lower bound on $s_*$ beyond this
formula is asserted.

\begin{theorem}[Uniform higher integrability of the unprotected density]
\label{thm:v54-uniform-integrability}
For every $h>0$ and $m\ge2$,
\begin{equation}\label{eq:v54-uniform-integrability}
 \sup_{R,n,k,t}\int_t^{t+h}
                 |m^2p_{n,R}(k,m,u)|^{q_*}\,du\le C_*(1+h).
\end{equation}
More precisely, for all $L\ge L_0$,
\begin{equation}\label{eq:v54-physical-height-tail}
 \sup_{R,n,k,t}\int_{[t,t+h]\cap\{P>L\}}P(u)\,du
 \le C(1+h)(L^{-1/144}+e^{-\kappa m}).
\end{equation}
The density is the complete original source, with no protection or
index allowance in its definition.
\end{theorem}
\begin{proof}
On $[t,t+h]$ apply Lemma~\ref{lem:v54-tail-principle} to the exact
split \eqref{eq:v54-positive-guard-split}, with $K=9$ and
$\alpha=1/16$. The needed protected height is
Proposition~\ref{prop:v54-protected-height}; its positive remainder
mass is \eqref{eq:v54-all-margin-window}. The total local mass is
bounded by \eqref{eq:v54-unmarked-window}. The complete exponential
cap and the strict inequalities for $s_*$ were established just above.
All interval-dependent constants are at most a fixed multiple of
$1+h$. This proves both assertions uniformly in the stated parameters.
\end{proof}

\paragraph{The role of physical boundaries.}
The proof uses a flow only on the fully protected source. Histories
with incidence, clearance or section defects enter the positive
remainder before any upper comparison. Their local mass, the polynomial
protected height and the exponential full-source cap are three
different inputs. The new bound does not infer height from mass and
does not remove a collision at grazing. It rules out lack of uniform
integrability of $P\log_+P$, but not unbounded thin density spikes.
